\section{生成函数、矩母函数与特征函数}
\label{sec:05-Probability/07-Transforms-and-Bounds/01}

一个推荐位一天的点击来自七个投放渠道，各渠道当天的点击数分布你都估出来了，彼此可以当作独立。总点击数落在某个区间的概率是多少？

按定义走，两个变量之和的分布是卷积：

\[
P(X+Y=k)=\sum_{j} P(X=j)\,P(Y=k-j).
\]

这在\hyperref[sec:05-Probability/05-Joint-Distributions/03]{``随机变量之和与卷积''}里已经推过。麻烦不在于它难懂，而在于它不肯化简：每算一个 \(k\) 都要跑一遍求和，七个渠道就是六层嵌套的卷积，中间结果的支撑还越滚越宽。

对数处理乘法时用的招数值得借鉴——\(\log(ab)=\log a+\log b\)，换一个坐标系，难做的运算就退化成好做的运算。概率论要找的是反方向的同一件事：一种编码方式，使得随机变量的相加对应编码后的相乘。

\subsection{同一个想法的三种实现}

三种编码的写法只差一个指数的形态：

\[
G_X(z)=\E[z^X],\qquad
M_X(t)=\E[e^{tX}],\qquad
\varphi_X(t)=\E[e^{itX}].
\]

它们分别叫概率母函数（PGF）、矩母函数（MGF）和特征函数。三者共享同一条性质，也正是引入它们的理由：若 \(X,Y\) 独立，则

\[
\E[f(X)f(Y)]=\E[f(X)]\,\E[f(Y)],
\]

而 \(z^{X+Y}=z^Xz^Y\)、\(e^{t(X+Y)}=e^{tX}e^{tY}\) 恰好把``和''拆成了乘积形式。于是

\[
G_{X+Y}=G_XG_Y,\qquad M_{X+Y}=M_XM_Y,\qquad \varphi_{X+Y}=\varphi_X\varphi_Y.
\]

\begin{center}
\begin{tikzpicture}[x=1cm,y=1cm,font=\footnotesize,>=stealth]
  \foreach \i/\h in {0/0.25,1/0.55,2/0.35,3/0.15}
    {\fill[blue!30] (0.3*\i,0) rectangle (0.3*\i+0.22,\h);}
  \draw[black!55] (-0.1,0) -- (1.35,0);
  \node at (0.6,-0.4) {\(p_X\)};
  \node at (1.75,0.3) {\(*\)};
  \foreach \i/\h in {0/0.5,1/0.4,2/0.2,3/0.1}
    {\fill[blue!30] (2.2+0.3*\i,0) rectangle (2.2+0.3*\i+0.22,\h);}
  \draw[black!55] (2.1,0) -- (3.55,0);
  \node at (2.8,-0.4) {\(p_Y\)};
  \node at (3.95,0.3) {\(=\)};
  \foreach \i/\h in {0/0.12,1/0.3,2/0.5,3/0.42,4/0.28,5/0.14,6/0.06}
    {\fill[orange!45] (4.4+0.3*\i,0) rectangle (4.4+0.3*\i+0.22,\h);}
  \draw[black!55] (4.3,0) -- (6.5,0);
  \node at (5.4,-0.4) {\(p_{X+Y}\)};
  \node[right,black!70,align=left] at (6.8,0.3) {每一根柱子\\都要跑一遍求和};
  \draw[->,thick,black!60] (1.7,-0.75) -- (1.7,-1.55);
  \node[right,black!70] at (1.85,-1.15) {取 \(\E[z^X]\)};
  \draw[<-,thick,black!60] (5.4,-0.75) -- (5.4,-1.55);
  \node[right,black!70] at (5.55,-1.15) {读 \(z^k\) 的系数};
  \draw[black!45] (0,-3.3) rectangle (1.35,-2.1);
  \draw[blue!65,thick] plot[domain=0.05:1.3,samples=25]
    (\x,{-3.25+0.9*(\x/1.3)*(\x/1.3)});
  \node at (0.68,-3.7) {\(G_X(z)\)};
  \node at (1.75,-2.7) {\(\times\)};
  \draw[black!45] (2.2,-3.3) rectangle (3.55,-2.1);
  \draw[blue!65,thick] plot[domain=2.25:3.5,samples=25]
    (\x,{-3.25+0.9*((\x-2.2)/1.35)*((\x-2.2)/1.35)*((\x-2.2)/1.35)});
  \node at (2.88,-3.7) {\(G_Y(z)\)};
  \node at (3.95,-2.7) {\(=\)};
  \draw[black!45] (4.4,-3.3) rectangle (5.75,-2.1);
  \draw[orange!85!black,thick] plot[domain=4.45:5.7,samples=30]
    (\x,{-3.25+0.9*((\x-4.4)/1.35)*((\x-4.4)/1.35)*((\x-4.4)/1.35)*((\x-4.4)/1.35)*((\x-4.4)/1.35)});
  \node at (5.08,-3.7) {\(G_XG_Y\)};
  \node[right,black!70,align=left] at (6.1,-2.7) {两个函数逐点相乘，\\没有求和};
\end{tikzpicture}

\emph{上排是概率空间里的卷积，下排是变换域里的乘法。同一件事换了坐标：走上面那条路要为每个可能的和值重算一次求和，走下面那条路只是两条曲线逐点相乘，代价与支撑的宽度无关。}
\end{center}

这张对照图是本节唯一需要记住的图像。后面三小节只是在回答同一个问题的不同侧面：编码用哪种指数，付得起什么代价。

\subsection{概率母函数把整张概率表存成幂级数}

若 \(X\) 只取非负整数，\(G_X(z)=\sum_k P(X=k)z^k\) 就是把整张概率质量表塞进一个幂级数，\(z^k\) 的系数即 \(P(X=k)\)。想取回某个概率，读系数即可；想取矩，在 \(z=1\) 处求导，\(G_X'(1)=\E[X]\)，二阶导给出 \(\E[X(X-1)]\)，方差随手可拼。这些求导的细节不必背，它们只是幂级数的常规操作。

值得看的是乘法性质怎么把组合计数一次做完。Bernoulli\((p)\) 的母函数只有两项，\(1-p+pz\)；\(n\) 个独立同分布的 Bernoulli 相加，母函数是 \((1-p+pz)^n\)，按二项式定理展开，\(z^k\) 的系数正是 \(\binom nk p^k(1-p)^{n-k}\)。二项分布的概率质量函数不是被推导出来的，是被乘法自动吐出来的。

开头那七个渠道也一样。若第 \(j\) 个渠道贡献的点击数母函数是 \(G_j\)，总数的母函数就是 \(\prod_{j=1}^{7}G_j(z)\)，展开后 \(z^k\) 的系数给出恰好 \(k\) 次点击的概率。真正需要盯住的是那句``可以当作独立''：一次大促、一次投放系统故障会让几个渠道同时异常，这时把母函数乘起来得到的分布已经不是总点击数的分布了。变换降低的是运算代价，它检验不了独立性假设。

母函数的辖区很清楚：非负整数。取值连续或者可以为负，这套幂级数就无从谈起。

\subsection{矩母函数强在矩，弱在存在性}

把 \(z^X\) 换成 \(e^{tX}\)，取值限制随即解除。\(M_X(t)=\E[e^{tX}]\) 的名字来自一个求导事实：对 \(t\) 求导会把 \(X\) 从指数里乘出来，于是 \(M_X^{(k)}(0)=\E[X^k]\)。逐阶求导的具体演算这里不展开，它对理解后文没有作用；要记住的是 \(M_X\) 在 \(t=0\) 附近的展开

\[
M_X(t)=1+t\E[X]+\tfrac{t^2}{2}\E[X^2]+\cdots
\]

把所有阶的矩打包在了一个函数里。下一节的 Markov 不等式只用一阶矩，本单元最后一节的 Chernoff 界之所以能给出指数衰减，靠的正是这个打包。

Gaussian 是矩母函数的样板。\(X\sim\mathcal N(\mu,\sigma^2)\) 时

\[
M_X(t)=\exp\Bigl(\mu t+\tfrac12\sigma^2t^2\Bigr),
\]

两个独立 Gaussian 相加，指数上的二次多项式相加，形式不变，只是 \(\mu\) 与 \(\sigma^2\) 各自累加。``独立 Gaussian 之和仍是 Gaussian''这条结论，在变换域里只是一次多项式加法。

但这里有一个硬条件，不能一掠而过：\(M_X(t)\) 未必存在。重尾分布常常对每一个 \(t>0\) 都有 \(\E[e^{tX}]=\infty\)。典型病例是对数正态分布，即 \(X=e^Y\) 而 \(Y\) 服从 Gaussian。此时 \(\E[e^{tX}]=\E[e^{te^Y}]\)，指数套着指数，而 Gaussian 密度的 \(e^{-y^2/(2\sigma^2)}\) 追不上 \(e^{te^y}\) 的爆炸，积分发散。有意思的是，这个分布的每一阶矩都有限——矩全存在，矩母函数照样不存在，两件事不能互推。关于重尾如何逐级瓦解均值、方差和指数矩，见\hyperref[sec:05-Probability/06-Common-Distributions/07]{``重尾分布：当均值与方差失效时''}。

\subsection{特征函数换来无条件的存在性}

把实指数换成虚指数，存在性问题就消失了：

\[
\varphi_X(t)=\E[e^{itX}],\qquad i^2=-1.
\]

因为 \(\abs{e^{itX}}=1\) 恒成立，这个期望的模不超过 \(1\)，对任何分布、任何实数 \(t\) 都有限。不需要矩条件，不需要尾部假设，重尾分布也照收不误。乘法性质完好保留：\(\varphi_{X+Y}=\varphi_X\varphi_Y\)。

代价在哪里？付在了两处。一是值域跑到了复平面，直觉不如实函数好用；二是矩不再免费——\(\varphi_X\) 处处存在，不等于它处处可导，\(k\) 阶导数在原点给出 \(i^k\E[X^k]\) 的前提仍然是 \(\E\abs{X}^k<\infty\)。存在性与光滑性是两笔账。

这个差别决定了一件具体的事：中心极限定理的经典证明为什么走特征函数而不走矩母函数。定理的结论要对一切有限方差的分布成立，其中包括大量矩母函数不存在的分布；证明的骨架是把``和的分布收敛''翻译成``特征函数逐点收敛''，再由反演唯一性翻译回去。若用矩母函数，第一步就要先假设它存在，定理的适用范围立刻缩水。完整论证见\hyperref[sec:05-Probability/08-Limit-Theorems/04]{``特征函数怎样证明中心极限定理''}。

三种变换其实是同一棵树的三个杈，只要相应表达式有意义，就有 \(M_X(t)=G_X(e^t)\)、\(\varphi_X(t)=G_X(e^{it})\)。选哪一个取决于你要什么：读系数或做递推用母函数；要矩、要指数型尾界，且零点附近矩母函数确实存在，用矩母函数；要对一般分布做严格的收敛论证，用特征函数。

\subsection{累积量：把乘积再压成加法}

矩母函数存在时，对它取一次对数，

\[
K_X(t)=\log M_X(t),
\]

独立和在这一层实现彻底的线性：\(K_{X+Y}=K_X+K_Y\)。前两阶累积量恰好是均值和方差，更高阶的累积量刻画偏斜与尾厚这些方差看不见的形状信息。

Gaussian 在这里露出了它的代数身份：\(K_X(t)=\mu t+\frac12\sigma^2t^2\)，三阶及以上累积量全为零。反过来也成立，三阶以上累积量全为零的分布只能是 Gaussian。独立 Gaussian 相加之后 \(K\) 仍是二次多项式，所以结果仍是 Gaussian——这句话与前面那句``指数上的多项式相加''是同一件事的两种说法。

\(K_X\) 还会在下下节再次登场。Chernoff 界的最优指数正是 \(\sup_t\,[ta-K_X(t)]\)，一个凸共轭的形状；独立和让 \(K\) 与 \(n\) 成正比，最优指数也随之与 \(n\) 成正比，指数衰减就是这样出来的。

到此为止，我们做的都是``把分布算清楚''这类工作，只是换了个更省力的坐标系。下一节把要求降下来：完整分布不知道，只攥着一个均值、一个方差，甚至只有几个事件的边缘概率，还能对坏事件说出什么。
